\scanpage{27}
\mspage{130}\textbf{--- Série A}\hfill\textbf{C. R. Acad. Sc. Paris, t. 279 (22 juillet 1974)}\par\noindent\rule{\linewidth}{.35pt}\par

point fixe et qui inverse l'orientation. Un morphisme entre \((\widetilde{\mathrm{T}},\theta)\) et \((\widetilde{\mathrm{T}}',\theta')\) est une classe d'isotopie \(\bar h\) d'homéomorphisme \(h\) de voisinages tubulaires de \(\widetilde{\mathrm{T}}\) sur \(\widetilde{\mathrm{T}}'\) tel que \(h\,\theta=\theta'\,h\).

9.2. On définit un foncteur \(\mathrm{Q}\) de \(\widetilde{\mathrm{TUBO}}\) dans TUB qui à un objet \((\widetilde{\mathrm{T}},\theta)\) de \(\widetilde{\mathrm{TUBO}}\) associe le voisinage tubulaire quotient de \(\widetilde{\mathrm{T}}\) par \(\theta\) et qui à un morphisme associe le morphisme quotient.

9.3. \textsc{Proposition.} --- \textit{\(\mathrm{Q}\) est une équivalence entre les catégories \(\widetilde{\mathit{TUBO}}\) et TUB.}

10. \textsc{La catégorie isotopique des \(\mathbf{Z}_2\)-1-complexes circularisés.}

10.1. C'est la catégorie \(\widetilde{\mathrm{K}}\,\omega_{\mathrm{I}}\) dont les objets sont les triplés \((\widetilde{\mathrm{K}},\widetilde\omega,\theta)\) où \((\widetilde{\mathrm{K}},\widetilde\omega)\) est un objet de \(\mathrm{K}\,\omega_{\mathrm{I}}\) et \(\theta\) un homéomorphisme involutif sans point fixe de \(\widetilde{\mathrm{K}}\) qui inverse la circularisation (on dira que \(\widetilde\omega\) est \textit{antisymétrique}). Un morphisme entre \((\widetilde{\mathrm{K}},\widetilde\omega,\theta)\) et \((\widetilde{\mathrm{K}}',\widetilde\omega',\theta')\) est une classe d'isotopie \(\bar h\) d'homéomorphisme \(h\) de \(\widetilde{\mathrm{K}}\) sur \(\widetilde{\mathrm{K}}'\) qui respecte les circularisations et tel que \(h\,\theta=\theta'\,h\).

10.2. Le foncteur \(\mathrm{T}\) induit un foncteur \(\widetilde{\mathrm{T}}\) de \(\widetilde{\mathrm{K}}\,\omega_{\mathrm{I}}\) dans \(\widetilde{\mathrm{TUBO}}\).

\textsc{Théorème.} --- \textit{\(\widetilde{\mathrm{T}}\) est une équivalence de catégories entre \(\widetilde{\mathrm{K}}\,\omega_{\mathrm{I}}\) et \(\widetilde{\mathit{TUBO}}\).}

\textsc{Corollaire.} --- \textit{Les catégories \(\widetilde{\mathrm{K}}\,\omega_{\mathrm{I}}\) et TUB sont équivalentes.}

11. \textsc{La catégorie isotopique des 1-complexes cocircularisés.}

11.1. Une \textit{cochaîne} sur \(\mathrm{K}\) est une application \(\gamma\) de \(\mathrm{A}\) (arcs de \(\mathrm{K}\)) dans \(\mathbf{Z}_2\) (i. e. \(\mathbf{Z}/2\,\mathbf{Z}\)) telle que \(\gamma\sigma=\gamma\) [\textit{voir} (\textsuperscript{1})]. Si \(\omega\) et \(\omega'\) sont deux circularisations de \(\mathrm{K}\) telles que pour tout sommet \(s\) de degré \(\geqq3\) (le degré \(\mathring{d}s\) de \(s\) est le nombre d'arcs issus de \(s\)) \(\omega_s\) et \(\omega'_s\) soient égaux ou inverses l'un de l'autre (ce qu'on notera \(\lvert\omega_s\rvert=\lvert\omega'_s\rvert\)), on désigne par \(\omega'_s/\omega_s\) leur « rapport » qui vaut \(1\) dans le premier cas, \(-1\) dans le second. Si on note \(\varphi\) l'isomorphisme du groupe multiplicatif \(\{-1,1\}\) dans \(\mathbf{Z}_2\), on définit un « cobord » \(\chi_{\omega\omega'}\), application définie sur l'ensemble des arcs qui n'ont pas d'extrémités de degré 1, qui à un arc d'extrémités \(s_0\) et \(s_1\) associe \(\varphi\,(\omega'_{s_1}/\omega_{s_1})-\varphi\,(\omega'_{s_0}/\omega'_{s_0})\), et qui à une boucle isolée associe \(0\).

11.2. \(\mathrm{K}\) étant donné, on définit une relation d'équivalence sur les couples \((\gamma,\omega)\) où \(\gamma\) est une cochaîne, \(\omega\) une circularisation par \((\gamma,\omega)\sim(\gamma',\omega')\Leftrightarrow\) pour tout sommet \(s\) de \(\mathrm{K}\) tel que \(\mathring d s\geqq3\) on a \(\lvert\omega'_s\rvert=\lvert\omega_s\rvert\) et \(\chi_{\omega\omega'}=\gamma'-\gamma\) (sur le domaine où \(\chi_{\omega\omega'}\) est défini). Une classe \(\eta\) de couples \((\gamma,\omega)\) sera appelée \textit{cocircularisation} de \(\mathrm{K}\).

11.3. \textsc{La catégorie isotopique des 1-complexes cocircularisés} est la catégorie \(\mathrm{K}\,\eta{\scriptstyle\mathrm{I}}\) dont les objets sont les couples \((\mathrm{K},\eta)\) où \(\eta\) est une cocircularisation du 1-complexe \(\mathrm{K}\). Un morphisme de \((\mathrm{K},\eta)\) dans \((\mathrm{K}',\eta')\) est la donnée d'une classe d'isotopie d'homéomorphisme de \(\mathrm{K}\) dans \(\mathrm{K}'\) qui respecte les cocircularisations et, pour chaque point isolé, chaque arc isolé, chaque boucle isolée \(b\) telle que \(\gamma(b)=0\), la donnée d'un signe \(+1\) ou \(-1\) (intuitivement, le signe indique si l'extension du morphisme au voisinage tubulaire conserve ou inverse l'orientation).

11.4. \textsc{Théorème.} --- \textit{Les catégories \(\mathrm{K}\,\eta_{\mathrm{I}}\) et TUB sont équivalentes.}

[On définit un foncteur de \(\mathrm{K}\,\eta_{\mathrm{I}}\) dans TUB qui à \((\mathrm{K},\eta)\) associe un « voisinage tubulaire standard » \(\mathrm{T}(\mathrm{K},\eta)\) et un foncteur quasi-inverse du précédent qui à un voisinage tubulaire de \(\mathrm{K}\) associe le couple \((\mathrm{K},\eta)\) où \(\eta\) est la « cocircularisation induite » sur \(\mathrm{K}\).]

11.5. \textsc{Terminologie.} --- Un \textit{circuit} de \(\mathrm{K}\) est un ensemble non vide d'arcs ordonné circulairement, l'extrémité d'un arc coïncidant avec l'origine du suivant. Une \textit{boucle} est le couple formé par un circuit et son inverse. Si \((\mathrm{K},\eta)\) est un 1-complexe cocircularisé
% archive sheet number: 26

